\section{来源与复现}
\label{chapter:S3}

% block-id: S3.report_repro.heading
\hypertarget{S3.report_repro.heading}{}
\subsection{数据身份与科学状态}

% block-id: S3.report_repro.p1
\hypertarget{S3.report_repro.p1}{}
各路线的数据清单以模型、参数、源码与运行身份连接原始 WXF、检查点、Ansys 工程、日志和后处理结果。JSON/CSV 保存复数实部与虚部、单位及股线身份；源字段存在单位歧义时，规范量值依据源模型定义解释，并记录转换依据。逐图映射进一步关联数据、绘图表达式、独立 PDF/SVG/PNG 及组图布局。清单分别说明标量积分、空间场和部分恢复结果的覆盖范围。

% block-id: S3.report_repro.p2
\hypertarget{S3.report_repro.p2}{}
失败案例保存触发输入、实际错误、针对性修订、停止条件和继续研究所需条件。几何验证的包装器退出码及内核退出信息缺失，RT0 q=18 原摘要的缺失指标和失败状态，以及较大 PEEC 规模的异常退出，均按原运行记录保留；纵向 PEEC 的 1.28 m 档状态为未运行。后处理验收使用独立身份，读者可据此区分数据恢复、数值求解和物理验证的状态。

% block-id: S3.report_repro.p3
\hypertarget{S3.report_repro.p3}{}
LaTeX 主稿维护正文、公式和图注，生成的 Markdown 提供阅读稿。数据集清单定位保留的数值及几何证据，科研计算通过文本教程与 CLI 执行。迁移核验包括相对链接、附件读取、文件 SHA256 和含空格目录中的轻量运行；文档构建记录给出线程、内存与环境信息，便于在具备相应依赖的设备上复现。

% block-id: S3.workflow.heading
\hypertarget{S3.workflow.heading}{}
\subsection{计算与文档构建}

% block-id: S3.workflow.p1
\hypertarget{S3.workflow.p1}{}
研究计算入口读取模型输入和求解配置，输出运行记录及相应检查结果；文档构建入口读取已保存量值、图件与 LaTeX 主稿，生成 Markdown 阅读稿并编译预览 PDF。输入、源码及输出哈希共同定义一次构建，所用单图保留矢量曲线和字体信息，组图记录面板身份与布局参数。下列来源表用于从正文观测量追溯到具体文件及其验收记录。

% block-id: S3.sources
\hypertarget{S3.sources}{}
来源清单：metrics.json; migration-receipt.json; reading-receipt.json.
\begin{table*}[!htbp]
\centering
\caption{冻结来源与完整性}
\label{tab:S3.sources-1}
\begingroup\footnotesize
\begin{tabularx}{\textwidth}{>{\raggedright\arraybackslash}X>{\raggedright\arraybackslash}X}
\toprule
来源 & SHA256\\
\midrule
docs/总控/当前运行索引.json & \seqsplit{3324976d4af8f58e60820a88ae978246443f78fbdd0cfa1050e86809110ae60d}\\
outputs/prewrite\_\allowbreak closure\_\allowbreak state.json & \seqsplit{80a66afa79c321a89f730b102af0e70125f6e1655372717d6b59f470c6520e71}\\
\bottomrule
\end{tabularx}
\endgroup
\end{table*}
